\section{Quasar finality binding}
\label{sec:quasar-finality-binding}

Lux finality is post-quantum: BLS aggregates (Beam, classical fast
path) plus ML-DSA per-validator attestations plus Pulsar threshold
Pulses (lattice / Module-LWE), composed by Prism into a HorizonCertificate
(LP-105). FHE results enter this finality through a
$\mathsf{nebula\_fhe\_root}$ commitment in the QuasarRoundDescriptor.

\subsection{The commitment}

\[
  \mathsf{nebula\_fhe\_root}
  \;=\;
  H\!\bigl(
    \mathsf{bootstrap\_key\_root}
    \;\|\;
    \mathsf{key\_switch\_key\_root}
    \;\|\;
    \mathsf{public\_param\_root}
    \;\|\;
    \mathsf{active\_circuit\_dag\_root}
    \;\|\;
    \mathsf{finalized\_evaluations\_root} \bigr)
\]

\paragraph{Field origins.}
$\mathsf{bootstrap\_key\_root}$ and $\mathsf{key\_switch\_key\_root}$ are
produced by the M-Chain MPC keygen ceremony (LP-019, LP-076). The
Lux FHE-Go runtime computes them deterministically from the threshold
share output. $\mathsf{public\_param\_root}$ commits to the
ParameterSetID lineage. $\mathsf{active\_circuit\_dag\_root}$ is the
merkle root of currently-installed FHE circuits.
$\mathsf{finalized\_evaluations\_root}$ accumulates Nebula-FHE
$\mathsf{FheReceipt}$ commitments produced during the round.

\paragraph{Hash suite.}
Production builds use SHA3-256 as the canonical hash (LP-105
hash-suite invariance). The rc1 scaffold ships a self-contained
SHA-256 to reduce the dependency surface during scaffold validation;
the production handover is a one-line swap to
\texttt{luxcpp/crypto/sha256}. Both runtimes carry the
\texttt{HashSuiteID} field in the round descriptor so a verifier can
unambiguously distinguish.

\subsection{F-Chain alias}

In F-Chain deployments (LP-013), $\mathsf{nebula\_fhe\_root}$ and
$\mathsf{fchain\_fhe\_root}$ are the same 32-byte field. LP-013 names
it from the F-Chain perspective; this paper names it from the
dual-runtime perspective. Specifications and code paths must use one
name consistently within a given module --- inside
\texttt{github.com/luxfi/consensus} the field is named
\texttt{nebula\_fhe\_root}; inside \texttt{github.com/luxfi/fchain}
it is named \texttt{fchain\_fhe\_root}.

\subsection{Soundness}

Theorem~\ref{thm:fhe-nebula-binding} (proofs/fhe/nebula-binding.tex)
formalizes the binding soundness: a forged binding reduces to either
a hash collision in $H$ or a Pulsar SUF-CMA forgery
(\texttt{proofs/pulsar/unforgeability.tex},
Theorem~\ref{thm:pulsar-suf-cma}). The plaintext content of the FHE
evaluation is private throughout --- the GPU only handles ciphertexts
and threshold key holders are TEE-protected (LP-065). What Quasar
finalizes is the \emph{commitment} to the encrypted result, not the
result itself. Decrypting the result requires the threshold MPC
ceremony; finality of the encrypted-state machine does not.

\subsection{Replay protection}

Every contributing sub-root is keyed by $\mathsf{KeyEraID} +
\mathsf{Generation}$, and the $\mathsf{HashSuiteID}$ is bound into the
transcript. Cross-epoch replay of an FHE result is rejected because
the resulting $\mathsf{nebula\_fhe\_root}$ does not match the
descriptor in the new round; Prism's transcript-equality check
(LP-105 \S Prism) fails. The proof appears in
\texttt{proofs/fhe/nebula-binding.tex}, Remark~\ref{rem:fhe-replay-protection}.
